\documentclass[12pt]{article}
\usepackage{amsmath, amssymb}
\usepackage[margin=1in]{geometry}
\setlength{\parskip}{6pt}
\title{QUBO Formulation: K-Medoids Prototype Selection}
\date{}
\begin{document}
\maketitle

\subsection*{K-Medoids Prototype Selection}

\paragraph{Problem Statement.}
Given a set of $N$ objects indexed by $V = \{0, 1, \dots, N-1\}$ and a pairwise distance matrix $D = (d_{ij})_{i,j \in V}$, the task is to select exactly $K$ objects to serve as medoids. Each object must be assigned to its nearest selected medoid. The objective is to minimize the total assignment distance, defined as the sum over all objects of their distance to the closest selected medoid.

\paragraph{Decision Variables.}
For each candidate object $i \in V$, we introduce a binary decision variable $x_i \in \{0, 1\}$. The variable $x_i = 1$ indicates that object $i$ is selected as a medoid, while $x_i = 0$ indicates it is not selected.

\paragraph{QUBO Derivation.}
\textbf{Step 1 --- Objective Function.} The original objective minimizes the total assignment distance. For QUBO encoding, we represent the cost contribution of each candidate medoid using precomputed coefficients $v_i \geq 0$, which capture the assignment cost associated with selecting medoid $i$. The objective is thus formulated as:
\begin{align}
    \min_{x \in \{0,1\}^N} \sum_{i \in V} v_i x_i.
\end{align}

\textbf{Step 2 --- Cardinality Constraint.} We enforce that exactly $K$ medoids are selected via the equality constraint:
\begin{align}
    \sum_{i \in V} x_i = K.
\end{align}

\textbf{Step 3 --- Quadratic Penalty Expansion.} We convert the equality constraint into a quadratic penalty term with weight $\lambda > 0$:
\begin{align}
    P(x) = \lambda \left( \sum_{i \in V} x_i - K \right)^2.
\end{align}
Expanding the square and applying the idempotent property $x_i^2 = x_i$ for binary variables yields:
\begin{align}
    P(x) &= \lambda \left( \sum_{i \in V} x_i^2 - 2K \sum_{i \in V} x_i + K^2 + \sum_{i \in V} \sum_{j \in V, j \neq i} x_i x_j \right) \\
    &= \lambda \left( \sum_{i \in V} x_i - 2K \sum_{i \in V} x_i + K^2 + 2 \sum_{i < j} x_i x_j \right) \\
    &= \lambda (1 - 2K) \sum_{i \in V} x_i + 2\lambda \sum_{i < j} x_i x_j + \lambda K^2.
\end{align}

\textbf{Step 4 --- Combined Hamiltonian.} Adding the objective and penalty terms gives the total QUBO Hamiltonian $H(x)$:
\begin{align}
    H(x) = \sum_{i \in V} v_i x_i + \lambda (1 - 2K) \sum_{i \in V} x_i + 2\lambda \sum_{i < j} x_i x_j + \lambda K^2.
\end{align}

\textbf{Step 5 --- QUBO Matrix Entries.} The Hamiltonian is expressed in the standard QUBO form $H(x) = \sum_{i \in V} Q_{ii} x_i + \sum_{i < j} Q_{ij} x_i x_j + c$. By collecting coefficients, we obtain the general closed-form expressions for the upper-triangular Q matrix entries and the constant offset:
\begin{align}
    Q_{ii} &= v_i + \lambda (1 - 2K), \quad \forall i \in V, \\
    Q_{ij} &= 2\lambda, \quad \forall i, j \in V \text{ with } i < j, \\
    c &= \lambda K^2.
\end{align}

\paragraph{Penalty Weight Justification.}
Let $V_{\max} = \max_{i \in V} |v_i|$ denote the maximum absolute cost coefficient in the objective. The maximum possible reduction in the objective function achievable by any single variable flip is bounded by $V_{\max}$. To guarantee that any violation of the cardinality constraint incurs an energy cost strictly greater than the maximum possible objective improvement, we require $\lambda > V_{\max}$. Under this condition, the penalty term dominates the objective landscape, ensuring that infeasible configurations are energetically unfavorable.

\paragraph{Correctness Argument.}
Minimizing $H(x)$ over the Boolean hypercube $\{0,1\}^N$ recovers the optimal medoid selection. First, any feasible solution satisfying $\sum_{i \in V} x_i = K$ incurs zero penalty, so $H(x)$ evaluates exactly to the objective function. Second, any infeasible solution violates the cardinality constraint by at least one unit, incurring a penalty of at least $\lambda$. Since $\lambda > V_{\max}$, the energy increase from constraint violation strictly exceeds any potential decrease in the objective function. Consequently, the global minimum of $H(x)$ must correspond to a feasible assignment with the minimum total assignment distance.

\paragraph{Illustrative Example.}
Consider a small instance with $N = 3$ objects and $K = 2$ required medoids. Let the precomputed cost coefficients be $v_0 = 10$, $v_1 = 12$, and $v_2 = 11$. We choose a penalty weight $\lambda = 15$, which satisfies $\lambda > \max\{10, 12, 11\}$. Substituting these values into the general formulas yields the concrete components:

Concrete objective: $10x_0 + 12x_1 + 11x_2$.

Concrete penalty terms: $15(x_0 + x_1 + x_2 - 2)^2 = -45(x_0 + x_1 + x_2) + 30(x_0 x_1 + x_0 x_2 + x_1 x_2) + 60$.

Resulting Hamiltonian:
\begin{align}
    H(x) &= -35x_0 - 33x_1 - 34x_2 + 30(x_0 x_1 + x_0 x_2 + x_1 x_2) + 60.
\end{align}

Enumerating all feasible configurations (exactly two variables set to 1):
\begin{align}
    x = (1, 1, 0) &\implies H = -35 - 33 + 30(1) + 60 = 22, \\
    x = (1, 0, 1) &\implies H = -35 - 34 + 30(1) + 60 = 21, \\
    x = (0, 1, 1) &\implies H = -33 - 34 + 30(1) + 60 = 23.
\end{align}
The global minimum occurs at $x^* = (1, 0, 1)$ with energy $H(x^*) = 21$, corresponding to selecting objects 0 and 2 as medoids.

\end{document}